\subsection{The Commutant and Bicommutant of a Module}

Let E be a ring and B a subset of E. The commutant (or centralizer) of B in E is the subring \(B'\) of E consisting of the elements that commute with every element of B. The commutant \(B''\) of \(B'\) is called the bicommutant (or bicentralizer) of B. We have \(B \subset B''\) , and \(B'\) coincides with its bicommutant (III, §1, No.~2, p.~429). The center of a subring B of E is \(B \cap B'\) ; the common center of \(B'\) and \(B''\) is \(B' \cap B''\) . If B is a commutative subring of E, then we have \(B \subset B'\) , and \(B''\) is the center of \(B'\) (III, §1, No.~2, p.~430).

Let A be a ring and M a left (resp. right) A-module. Let us apply these definitions to the case when E is the endomorphism ring of the additive group of M and B the ring of homotheties \(A_{M}\) of M. The commutant \(A_{M}^{\prime}\) of \(A_{M}\) in E is called the commutant of M; it is the endomorphism ring of the A-module M. The bicommutant \(A_{M}^{\prime\prime}\) of \(A_{M}\) in E is called the bicommutant of M; it is the endomorphism ring of the countermodule of M (VIII, p.~8, Definition 3). We have \(A_{M} \subset A_{M}^{\prime\prime}\) , the ring \(A_{M} \cap A_{M}^{\prime}\) is the center of \(A_{M}\) , and \(A_{M}^{\prime} \cap A_{M}^{\prime\prime}\) is the common center of \(A_{M}^{\prime}\) and \(A_{M}^{\prime\prime}\) .

DEFINITION 1. --- The A-module M is balanced if we have \(\mathrm{A}_{\mathrm{M}} = \mathrm{A}_{\mathrm{M}}^{\prime \prime}\) .

If the A-module M is balanced, then the rings \(A_{M}\) and \(A_{M}^{\prime}\) have the same center \(A_{M} \cap A_{M}^{\prime}\) . The module M is a balanced A-module if and only if it is a balanced \(A_{M}\) -module. The countermodule of an A-module M is faithful and balanced.

When the ring A is commutative, the bicommutant \(A_{M}^{\prime\prime}\) of M is the center of \(\mathrm{A}_{\mathrm{M}}^{\prime}=\mathrm{End}_{\mathrm{A}}(\mathrm{M})\) ; that M is balanced means that the center of \(\mathrm{End}_{\mathrm{A}}(\mathrm{M})\) is reduced to the homotheties.

For any element \(a\) of A, denote by \(\delta_{a}\) the right homothety \(x\mapsto xa\) from A to A and by \(\gamma_{a}\) the left homothety \(x\mapsto ax\) (I, §8, No.~1, p.~97). The map \(a\mapsto \delta_{a}\) is a ring isomorphism from \(\mathrm{A}^{\circ}\) to the commutant of the left A-module \(\mathrm{A}_s\) (cf.~II, §10, No.~7, p.~349, applied to \(\mathrm{I} = \{1\}\) ). The map \(a\mapsto \gamma_{a}\) is a ring isomorphism from A to the commutant of the right A-module \(\mathrm{A}_d\) (loc. cit.). If we identify A with the commutant of \(\mathrm{A}_d\) through this map, then the countermodule of \(\mathrm{A}_d\) is identified with \(\mathrm{A}_s\) ; consequently, the A-module \(\mathrm{A}_d\) is balanced. Likewise, the A-module \(\mathrm{A}_s\) is balanced.

Let n be an integer \(\geqslant 1\) . We view \(A^{n}\) as a left \(\mathbf{M}_{n}(\mathrm{A})\) -module (loc. cit.). The mapping that sends \(m \in \mathbf{M}_{n}(\mathrm{A})\) to the endomorphism \(x \mapsto mx\) of the A-module \(A_{d}^{n}\) is a ring isomorphism from \(\mathbf{M}_{n}(\mathrm{A})\) to the commutant of the right A-module \(A_{d}^{n}\) (loc. cit.).

PROPOSITION 1. --- Let \((\mathrm{A}_{i})_{i\in\mathrm{I}}\) be a family of rings, and for every \(i\in I\) , let \(M_{i}\) be an \(A_{i}\) -module. Set \(A=\prod A_{i}, M=\prod M_{i}\) , and \(N=\bigoplus M_{i}\) . Endow M with the structure of an A-module with law of action \(((a_{i}),(x_{i}))\mapsto(a_{i}x_{i})\) . The set N is an A-submodule of M.

\begin{enumerate}
\def\labelenumi{\alph{enumi})}
\item
  The mapping \((u_i) \mapsto \prod u_i\) from \(\prod \mathrm{End}_{\mathbf{Z}}(\mathrm{M}_i)\) to \(\mathrm{End}_{\mathbf{Z}}(\mathrm{M})\) (II, §1, No.~5, p.~200) restricts to ring isomorphisms from \(\prod (\mathrm{A}_i)_{\mathrm{M}_i}\) , \(\prod (\mathrm{A}_i)'_{\mathrm{M}_i}\) , and \(\prod (\mathrm{A}_i)'''_{\mathrm{M}_i}\) to \(\mathrm{A}_{\mathrm{M}}\) , \(\mathrm{A}_{\mathrm{M}}'\) , and \(\mathrm{A}_{\mathrm{M}}''\) , respectively.
\item
  The mapping \((u_i) \mapsto \bigoplus u_i\) from \(\prod \operatorname{End}_{\mathbf{Z}}(\mathrm{M}_i)\) to \(\operatorname{End}_{\mathbf{Z}}(\mathrm{N})\) (II, §1, No.~6, p.~203) restricts to ring isomorphisms from \(\prod (\mathrm{A}_i)_{\mathrm{M}_i}\) , \(\prod (\mathrm{A}_i)'_{\mathrm{M}_i}\) , and \(\prod (\mathrm{A}_i)'''_{\mathrm{M}_i}\) to \(\mathrm{A_N}\) , \(\mathrm{A_N}'\) , and \(\mathrm{A_N}''\) , respectively.
\end{enumerate}

The ring isomorphisms defined in Proposition 1 are called canonical. The product \(\prod (\mathrm{A}_i)_{\mathrm{M}_i}\) is often identified with \(\mathrm{A_M}\) , and \(\prod (\mathrm{A}_i)'_{\mathrm{M}_i}\) with \(\mathrm{A_M}'\) , etc., through these isomorphisms.

The mapping \(\varphi : (u_i) \mapsto \prod u_i\) from \(\prod \operatorname{End}_{\mathbf{Z}}(\mathrm{M}_i)\) to \(\operatorname{End}_{\mathbf{Z}}(\mathrm{M})\) is an injective ring homomorphism. By the definition of the A-module structure on M, we have \(\varphi(\prod (\mathrm{A}_i)_{\mathrm{M}_i}) = \mathrm{A}_{\mathrm{M}}\) . Let \(u \in \mathrm{A}_{\mathrm{M}}'\) . For every \(i \in \mathrm{I}\) , denote by \(h_i\) the element of A with all components equal to 1 except that of index \(i\) , which is equal to 0. If \(x\) is an element of M with component of index \(i\) equal to 0, then we have \(x = h_i x\) , and therefore \(\operatorname{pr}_i(u(x)) = \operatorname{pr}_i(u(h_i x)) = \operatorname{pr}_i(h_i u(x)) = 0\) . Consequently, there exists a unique group homomorphism \(u_i: \mathrm{M}_i \mapsto \mathrm{M}_i\) such that \(\operatorname{pr}_i(u(y)) = u_i(\operatorname{pr}_i(y))\) for every \(y \in \mathrm{M}\) . We have \(u = \prod u_i\) . Since the mapping \(u\) is A-linear, the mapping \(u_i\) is \(\mathrm{A}_i\) -linear for every \(i \in \mathrm{I}\) . This proves that the image of \(\prod (\mathrm{A}_i)'_{\mathrm{M}_i}\) under \(\varphi\) contains \(\mathrm{A}_{\mathrm{M}}'\) ; the reverse inclusion is obvious. By applying this to the countermodule of M, we deduce that \(\varphi\) induces an isomorphism from \(A_{M}^{\prime \prime}\) to \(\prod_i(A_i)_{M_i}^{\prime \prime}\) . This proves assertion a).

The proof of b) is the same as that of a) mutatis mutandis.

PROPOSITION 2. --- Let A be a ring, and let M be an A-module. Let I be a set. Then the bicommutant of the A-module \(\mathrm{M}^{(I)}\) coincides with the ring of homotheties of the \(A_{M}^{\prime\prime}\) -module \(\mathrm{M}^{(I)}\) .

For \(i \in I\) , we denote the projection homomorphism from \(M^{(I)}\) to M by \(\pi_{i} : (x_{j})_{j \in I} \mapsto x_{i}\) and the corresponding canonical injection (I, §4, No.~9, p.~47) by \(\iota_{i} : M \to M^{(I)}\) .

Let \(u\) be an element of \(\mathrm{End}_{\mathrm{A}}(\mathrm{M}^{(\mathrm{I})})\) . For all \(i,j \in \mathrm{I}\) , the composition \(u_{i,j} = \pi_j \circ u \circ \iota_i\) belongs to the commutant \(\mathrm{A}_{\mathrm{M}}'\) of \(\mathrm{M}\) . For every element \(b\) of \(\mathrm{A}_{\mathrm{M}}''\) and every \((x_i) \in \mathrm{M}^{(\mathrm{I})}\) , we have the relations

\[
b u \left(\left(x _ {i}\right) _ {i \in \mathrm{I}}\right) = b \left(\sum_ {i \in \mathrm{I}} u _ {i, j} \left(x _ {i}\right)\right) _ {j \in \mathrm{I}} = \left(\sum_ {i \in \mathrm{I}} u _ {i, j} \left(b x _ {i}\right)\right) _ {j \in \mathrm{I}} = u \left(b \left(x _ {i}\right) _ {i \in \mathrm{I}}\right).
\]

The homothety \(b_{\mathrm{M}^{(\mathrm{I})}}\) therefore belongs to the bicommutant of the A-module \(\mathrm{M}^{(\mathrm{I})}\) .

Conversely, let b be an element of the bicommutant of \(\mathrm{M}^{(\mathrm{I})}\) . For \(i,j\in I\) , set \(b_{i,j}=\pi_{j}\circ b\circ\iota_{i}\) . Take \(i,j\in I\) with \(i\neq j\) . Since \(\iota_{j}\circ\pi_{j}\) belongs to the commutant of the A-module \(\mathrm{M}^{(\mathrm{I})}\) , we have

\[
b _ {i, j} = \pi_ {j} \circ \iota_ {j} \circ \pi_ {j} \circ b \circ \iota_ {i} = \pi_ {j} \circ b \circ \iota_ {j} \circ \pi_ {j} \circ \iota_ {i} = 0.
\]

Likewise, we have

\[
b _ {j, j} = \pi_ {j} \circ b \circ \iota_ {j} = \pi_ {i} \circ \iota_ {i} \circ \pi_ {j} \circ b \circ \iota_ {j} = \pi_ {i} \circ b \circ \iota_ {i} \circ \pi_ {j} \circ \iota_ {j} = b _ {i, i}.
\]

Moreover, \(b_{i,i}\) belongs to \(A_{M}^{\prime\prime}\) . It follows that b coincides with a homothety of the \(A_{M}^{\prime\prime}\) -module \(M^{(I)}\) .

